% Ljung-Box (TesteLjungBox do robô, 10 defasagens, 2 parâmetros ARMA, como na
% validação numérica): autocorrelação dos resíduos do ARIMA(1,1,1), dos seus
% quadrados e dos quadrados padronizados pelo GARCH; embaixo, como a
% estatística Q se acumula defasagem a defasagem até o valor crítico.
\documentclass[border=4pt]{standalone}
% Preâmbulo comum dos gráficos que explicam os métodos (30/10/2026).
% Dados e números: dados/resultados/metodos_ilustrados_2026-10-30/, gerados por
% ferramentas/MetodosIlustrados.mq5 com as funções do próprio robô. Nenhum
% número é digitado aqui: os do texto vêm de valores.tex (\vX para o texto,
% \pX para coordenada).
\usepackage[T1]{fontenc}
\usepackage{amsmath}
\usepackage{mathptmx}  % Times, como as figuras do TG
\usepackage{pgfplots}
\usepgfplotslibrary{groupplots,dateplot,colormaps,fillbetween}
\usetikzlibrary{arrows.meta,positioning}
\pgfplotsset{compat=1.18}
\definecolor{cmod}{RGB}{0,90,181}     % modelo / ajuste
\definecolor{cdes}{RGB}{220,50,32}    % destaque, limite, valor crítico
\definecolor{cdad}{RGB}{120,120,120}  % dado observado
\newcommand{\arq}[1]{../../dados/resultados/metodos_ilustrados_2026-10-30/#1}
\input{\arq{valores.tex}}
% Figuras 08 a 14 (o robô operando): traço por barra da refação de 05/10/2026,
% resumido por ferramentas/scripts/robo_em_acao.sh. Macros \vR e \pR.
\newcommand{\arqr}[1]{../../dados/resultados/robo_em_acao_2026-10-30/#1}
\input{\arqr{valores.tex}}
\definecolor{cari}{RGB}{230,140,0}   % ARIMA puro
\definecolor{cok}{RGB}{0,140,70}     % compra liberada
\usetikzlibrary{shapes.misc}
\newcommand{\passo}[1]{\tikz[baseline=(n.base)]\node[circle,fill=black!75,text=white,inner sep=1pt,font=\footnotesize\bfseries](n){#1};}
% eixo de datas comum às figuras do traço (jul/2019 a jun/2021)
\pgfplotsset{traco/.style={date coordinates in=x, xmin=2019-07-01, xmax=2021-06-30,
    xtick={2019-07-01,2020-01-01,2020-07-01,2021-01-01,2021-06-30},
    xticklabels={jul/2019,jan/2020,jul/2020,jan/2021,jun/2021}}}
\pgfplotsset{
  base/.style={table/col sep=comma, grid=major, grid style={gray!20},
    tick label style={font=\small}, label style={font=\small},
    title style={font=\small\bfseries, align=center},
    every axis plot/.append style={line width=0.9pt},
    /pgf/number format/use comma, /pgf/number format/1000 sep={.}},
}
% faixa vertical nos pregões do traço em que a condição #1 vale, na cor #2,
% do piso #3 ao teto #4 do eixo
\newcommand{\faixa}[4]{%
  \addplot[name path=topo, draw=none, const plot, forget plot] table[x=data, y expr={(#1) ? #4 : #3}]{\arqr{dia_a_dia.csv}};
  \addplot[name path=piso, draw=none, forget plot] coordinates {(2019-07-01,#3) (2021-06-30,#3)};
  \addplot[#2, forget plot] fill between[of=topo and piso];}

\begin{document}
\begin{tikzpicture}
\begin{groupplot}[base, group style={group size=3 by 2, horizontal sep=1.1cm, vertical sep=2.2cm,
    ylabels at=edge left},
  width=6.1cm, height=4.6cm]
% linha 1: autocorrelação
\nextgroupplot[title={Resíduos do ARIMA $\hat\varepsilon_t$}, ybar, bar width=3pt,
  xmin=0, xmax=21, ymin=-0.2, ymax=0.2, ylabel={Autocorrelação $\rho_k$}, xlabel={Defasagem $k$},
  /pgf/number format/fixed]
\addplot[fill=cdad, draw=none] table[x=k, y=residuo]{\arq{lb_acf.csv}};
\draw[cdes, dashed] (axis cs:0,\pLbBanda) -- (axis cs:21,\pLbBanda);
\draw[cdes, dashed] (axis cs:0,-\pLbBanda) -- (axis cs:21,-\pLbBanda);
\nextgroupplot[title={Quadrados $\hat\varepsilon_t^2$ (antes do GARCH)}, ybar, bar width=3pt,
  xmin=0, xmax=21, ymin=-0.2, ymax=0.2, xlabel={Defasagem $k$}, /pgf/number format/fixed]
\addplot[fill=cdad, draw=none] table[x=k, y=quadrado]{\arq{lb_acf.csv}};
\draw[cdes, dashed] (axis cs:0,\pLbBanda) -- (axis cs:21,\pLbBanda);
\draw[cdes, dashed] (axis cs:0,-\pLbBanda) -- (axis cs:21,-\pLbBanda);
\nextgroupplot[title={Padronizados $(\hat\varepsilon_t/\sigma_t)^2$ (depois)}, ybar, bar width=3pt,
  xmin=0, xmax=21, ymin=-0.2, ymax=0.2, xlabel={Defasagem $k$}, /pgf/number format/fixed]
\addplot[fill=cmod, draw=none] table[x=k, y=padronizado_quadrado]{\arq{lb_acf.csv}};
\draw[cdes, dashed] (axis cs:0,\pLbBanda) -- (axis cs:21,\pLbBanda);
\draw[cdes, dashed] (axis cs:0,-\pLbBanda) -- (axis cs:21,-\pLbBanda);
% linha 2: Q acumulado
\nextgroupplot[xmin=0.5, xmax=10.5, ymin=0, ymax=32, xtick={1,...,10}, ylabel={$Q$ acumulado},
  xlabel={Até a defasagem $k$}, title={$Q = \vLbQe$, $p = \vLbPe$:\\sobra autocorrelação}]
\addplot[cmod!60!black, mark=*, mark size=1.5pt] table[x=k, y=residuo]{\arq{lb_parcelas.csv}};
\draw[cdes, dashed, line width=1pt] (axis cs:0.5,\pLbCrit) -- (axis cs:10.5,\pLbCrit);
\nextgroupplot[xmin=0.5, xmax=10.5, ymin=0, ymax=32, xtick={1,...,10},
  xlabel={Até a defasagem $k$}, title={$Q = \vLbQeq$, $p = \vLbPeq$:\\a volatilidade se agrupa}]
\addplot[cmod!60!black, mark=*, mark size=1.5pt] table[x=k, y=quadrado]{\arq{lb_parcelas.csv}};
\draw[cdes, dashed, line width=1pt] (axis cs:0.5,\pLbCrit) -- (axis cs:10.5,\pLbCrit);
\nextgroupplot[xmin=0.5, xmax=10.5, ymin=0, ymax=32, xtick={1,...,10},
  xlabel={Até a defasagem $k$}, title={$Q = \vLbQzq$, $p = \vLbPzq$:\\o GARCH capturou o agrupamento}]
\addplot[cmod, mark=*, mark size=1.5pt] table[x=k, y=padronizado_quadrado]{\arq{lb_parcelas.csv}};
\draw[cdes, dashed, line width=1pt] (axis cs:0.5,\pLbCrit) -- (axis cs:10.5,\pLbCrit);
\node[cdes, font=\footnotesize, anchor=south east] at (axis cs:10.4,\pLbCrit) {valor crítico em $k = 10$: $\vLbCrit$};
\end{groupplot}
\node[font=\footnotesize, align=center, text=black!75] at ([yshift=-1.45cm]group c2r2.south)
  {Em cima, o tracejado é a faixa de ruído $\pm 1{,}96/\sqrt{n} = \pm\vLbBanda$.
   Embaixo, $Q = n(n+2)\sum_{k=1}^{10} \rho_k^2/(n-k)$ soma as barras de cima ao quadrado;\\
   passar do tracejado ($\chi^2$ com $\vLbGl$ graus de liberdade, 5\%) quer dizer que sobrou autocorrelação.};
\end{tikzpicture}
\end{document}
